\documentclass[11pt]{article}
\usepackage[margin=1in]{geometry}
\usepackage{amsmath,amssymb}
\usepackage{hyperref}

\title{Literature Status: Universal Differentiability Sets in Laakso Space}
\author{Automatic Math Discovery Run \texttt{fa\_banach\_001}}
\date{July 3, 2026}

\begin{document}
\maketitle

\section*{Status}

This is a lightweight literature-status packet. The source paper
\cite{Capolli2022} lists the existence of universal differentiability sets
in Laakso spaces as still open. A later paper \cite{EBPS2024} explicitly
answers that question affirmatively.

\section*{Source question}

In the final section, ``Further developments,'' Capolli's overview
\cite{Capolli2022} states that the existence of UDS in Laakso spaces is still
open. It explains that the earlier paper of Capolli--Pinamonti--Speight
proved only that the classical approach does not work in the Laakso space,
leaving the existence of UDS open.

\section*{Supporting answer}

Eriksson-Bique, Pinamonti, and Speight \cite{EBPS2024} explicitly identify
the same gap in their introduction: previous work left open whether
measure-zero UDS exist in Laakso space, and their paper answers this
question by a different method. Their Theorem 2 states that there is a Borel
set \(N\subset F\) with
\[
        \mathcal H^Q(N)=0
\]
such that every Lipschitz map \(f:F\to\mathbb R\) is differentiable at some
point of \(N\). Thus the Laakso-space UDS existence question from
\cite{Capolli2022} is answered affirmatively in the later literature.

\section*{Scope limitations}

This identification answers only the UDS existence item for the standard
Laakso space \(F\). It does not settle the other broad further-development
questions in \cite{Capolli2022}, such as Laakso-type spaces based on other
self-similar sets, non-Abelian Carnot models, or general inverse-limit
spaces.

\begin{thebibliography}{9}

\bibitem{Capolli2022}
M. Capolli,
\emph{An Overview on Laakso Spaces},
arXiv:2211.05681.

\bibitem{EBPS2024}
S. Eriksson-Bique, A. Pinamonti, and G. Speight,
\emph{Universal Differentiability Sets in Laakso Space},
arXiv:2406.07660.

\end{thebibliography}

\end{document}
